\section{Fundamentos de los modelos de lenguaje: partiendo del autocompletado}

\subsection{Qué es un modelo de lenguaje}

\begin{importantbox}{La esencia del modelo de lenguaje}
Un modelo de lenguaje (Language Model) es, en esencia, un sistema de ``autocompletado sofisticado'' (fancy autocomplete). Dado un fragmento de texto previo (stem), el modelo predice la palabra (token) siguiente más probable.
\end{importantbox}

El docente parte del ejemplo más intuitivo: dado ``It's raining cats and \_\_\_\_'', casi cualquier persona predice que la siguiente palabra es ``dogs''. Esto es precisamente lo que hace un modelo de lenguaje: predecir la siguiente palabra a partir del contexto.

Este proceso puede extenderse para predecir varias palabras seguidas. Por ejemplo, dado ``To be or not'', el modelo primero predice ``to'', lo concatena al final del contexto y después predice ``be'', con lo que se recupera la frase completa ``To be or not to be''. Esta forma de predecir un token a la vez y retroalimentar el resultado se llama \textbf{decodificación autorregresiva} (Autoregressive Decoding).

\begin{knowledgebox}{Decodificación autorregresiva (Autoregressive Decoding)}
La decodificación autorregresiva es el paradigma básico con el que los LLM actuales generan texto:
\begin{enumerate}
    \item Dada una secuencia de entrada $x_1, x_2, \dots, x_t$
    \item El modelo predice la distribución de probabilidad del siguiente token $P(x_{t+1} \mid x_1, \dots, x_t)$
    \item Se muestrea un token $x_{t+1}$ de esa distribución
    \item Se concatena $x_{t+1}$ al final de la secuencia y se repiten los pasos 2-3
    \item Hasta generar el token especial de fin de oración (end-of-sentence)
\end{enumerate}
\end{knowledgebox}

\subsection{Encajar el problema en la estructura de ``rellenar el espacio en blanco''}

El motivo por el que los modelos de lenguaje despertaron tanto interés es que muy diversos tipos de tareas pueden encajarse (embed) con ingenio en esta estructura de problema de ``rellenar el espacio en blanco'':

\begin{itemize}
    \item \textbf{Problemas matemáticos}: ``I have two apples and I eat one, I'm left with \_\_\_\_'' $\rightarrow$ ``one'' — el modelo hizo una operación aritmética
    \item \textbf{Razonamiento por analogía}: ``Paris is to France as Tokyo is to \_\_\_\_'' $\rightarrow$ ``Japan'' — el modelo hizo una analogía
    \item \textbf{Consulta de hechos}: ``Pizza was invented in \_\_\_\_'' $\rightarrow$ ``Naples, Italy'' — el modelo hizo una recuperación de conocimiento
\end{itemize}

El docente señala en particular que el razonamiento por analogía atormentó durante mucho tiempo a los investigadores de PLN. En la época de los embeddings de Word2Vec, los investigadores se esforzaban mucho y aun así les costaba lograr que el modelo resolviera bien las analogías, mientras que un estudiante de preparatoria las resuelve sin dificultad en el examen SAT. Este contraste tan marcado pone de relieve el salto de capacidad que trajeron los LLM.

\subsection{Construir un modelo de lenguaje estadístico}

El docente usa un modelo de lenguaje bayesiano (n-gram) para mostrar el método más básico de modelado del lenguaje. Este método se propuso ya en la década de 1980 y su idea central es el ``conteo sofisticado'' (fancy counting).

Con la célebre frase de Dickens ``It was the best of times, it was the worst of times'' como corpus:

\begin{enumerate}
    \item \textbf{Preprocesamiento del texto}: convertir a minúsculas, quitar la puntuación y agregar los tokens especiales de inicio (start-of-sentence) y fin de oración (end-of-sentence)
    \item \textbf{Construcción del diccionario de n-gramas}: contar cuántas veces aparecen todos los unigramas, bigramas, trigramas y 4-gramas
    \item \textbf{Estimación de la probabilidad condicional}: para el stem ``it was the'', contar las palabras que le siguen y su frecuencia
\end{enumerate}

Dado el stem ``it was the'', la distribución de probabilidad de las palabras siguientes en el corpus es:

\begin{table}[H]
\centering
\begin{tabular}{lcc}
\toprule
\textbf{Palabra siguiente} & \textbf{Número de apariciones} & \textbf{Probabilidad} \\
\midrule
age       & 2 & 1/3 \\
best      & 1 & 1/6 \\
epoch     & 2 & 1/3 \\
worst     & 1 & 1/6 \\
\bottomrule
\end{tabular}
\caption{Distribución de probabilidad condicional basada en estadísticas de n-gramas}
\end{table}

Al muestrear al azar de este diccionario de probabilidades, este modelo de lenguaje puede convertirse en un \textbf{modelo generativo}. Pero el resultado generado suele caer en un ``ciclo de probabilidad'' (probability loop):

\begin{quote}
\textit{``It was the best of times, it was the worst of times, it was the worst of times, it was the worst of times, it was the age of wisdom, it was the age of foolishness...''}
\end{quote}

\begin{warningbox}{La relación intuitiva entre el ciclo de probabilidad y la alucinación}
Este ejemplo tan sencillo muestra un fenómeno importante: el modelo cae en un ciclo repetitivo y produce sin cesar ``it was the worst of times''. Esto no se debe a que el modelo sea ``pesimista'', sino a que la ventana de contexto es demasiado pequeña, lo que hace que el modelo no sepa cuándo ``salir'' del patrón repetido. El mecanismo subyacente cuando escuchas hablar de que un LLM ``hallucinating'' (alucina) es similar: el modelo genera texto en una región singular de la distribución de probabilidad, no sabe qué decir y entonces produce al azar contenido que parece razonable pero que en realidad es incorrecto.
\end{warningbox}

\subsection{Resumen de la sección}

La tarea central de un modelo de lenguaje es la \textbf{predicción de la siguiente palabra} (next-word prediction). Desde los modelos estadísticos de n-gramas más simples hasta los LLM Transformer actuales, este paradigma básico no ha cambiado. Todo comportamiento que parece ``inteligente'' —hacer operaciones matemáticas, resolver analogías, consultar hechos— se logra, en esencia, encajando el problema en la estructura de ``rellenar el espacio en blanco''. Las limitaciones del modelo n-gram (el ciclo de probabilidad, la ventana de contexto pequeña) también anticipan los retos comunes que enfrentan los LLM.
